% Generated draft: requires mathematical and source review.
\documentclass[11pt]{article}
\usepackage[margin=25mm]{geometry}
\usepackage{fontspec}
\setmainfont{Latin Modern Roman}
\usepackage{amsmath,amssymb,mathtools,unicode-math}
\setmathfont{Latin Modern Math}
\usepackage{enumitem,needspace,xurl,hyperref}
\setlength{\parindent}{0pt}
\setlength{\parskip}{5pt}
\setlength{\emergencystretch}{3em}
\newcommand{\N}{\mathbb{N}}
\newcommand{\Z}{\mathbb{Z}}
\newcommand{\Q}{\mathbb{Q}}
\newcommand{\R}{\mathbb{R}}
\newcommand{\C}{\mathbb{C}}
\newcommand{\F}{\mathbb{F}}
\newcommand{\A}{\mathbb{A}}
\newcommand{\PP}{\mathbb{P}}
\newcommand{\E}{\mathbb{E}}
\newcommand{\eps}{\varepsilon}
\newcommand{\dd}{\,\mathrm{d}}
\begin{document}
\section*{UnsolvedMath ID: 30001725}
\textbf{Source problem code:} catalogue entry
\subsection{Definitions and mathematical statement}
\textbf{Research question.} Delimit the field of numerical tables and formulate criteria distinguishing a formal numerical table from a material array whose role is merely to organize numerical or nonnumerical data spatially.

Here, ``formal numerical table'' and ``material array'' designate the categories whose distinction is being sought; the supplied record does not define them mathematically. Its description of a material array specifies spatial organization of data, but supplies no necessary or sufficient conditions for membership in either category. The question concerns historical and conceptual classification. No universe of objects, equivalence relation, boundary conventions, or criteria for an adequate delimitation are specified. Consequently, a precise quantified mathematical proposition cannot be extracted without adding assumptions absent from the source. This draft preserves the research question and leaves those qualifications unresolved.
\subsection{Short English statement}
The task is to explain the scope of numerical tables and distinguish formal tables from spatial arrangements of numerical or nonnumerical data. The supplied question is historiographical and leaves the relevant definitions and classification criteria open.
\subsection{Sources}
\begin{itemize}
\item Supplied local UnsolvedMath registry record, permanent ID \(30001725\), title: \emph{Boundaries Between Formal Tables and Material Arrays}; exact source code: \texttt{catalogue entry}. Collection: Oberwolfach Reports Open Problems. The record gives 2011 as the proposed year and identifies no individual proposer.
\item \emph{Mini-Workshop: History of Numerical and Graphical Tables} (2011), Oberwolfach Reports, pp.~639--689, DOI: \texttt{10.4171/owr/2011/12}. Supplied source and literature reference, labeled ``Oberwolfach Report 12/2011'': \url{https://doi.org/10.4171/owr/2011/12}. The supplied background describes these reports as documenting mathematical workshops and their open-problem discussions.
\item Supplied literature assessment dated \texttt{2026-08-21}: the record labels the question open, describes the definition and delimitation of numerical tables as an ongoing historical research problem, and reports that no unique generally accepted resolution was located. This is supplied dated curation derived from the Oberwolfach Reports material, not an independently verified literature assessment; the status and cited claims require independent verification.
\item The registry assigns statement status \texttt{corrected\_verified} and states that the official report confirms the historiographical question and that an extraction label was removed. These are supplied source-verification claims, not verification performed for this draft.
\end{itemize}
\end{document}
